% Options for packages loaded elsewhere
\PassOptionsToPackage{unicode}{hyperref}
\PassOptionsToPackage{hyphens}{url}
\documentclass[
]{article}
\usepackage{xcolor}
\usepackage{amsmath,amssymb}
\setcounter{secnumdepth}{-\maxdimen} % remove section numbering
\usepackage{iftex}
\ifPDFTeX
  \usepackage[T1]{fontenc}
  \usepackage[utf8]{inputenc}
  \usepackage{textcomp} % provide euro and other symbols
\else % if luatex or xetex
  \usepackage{unicode-math} % this also loads fontspec
  \defaultfontfeatures{Scale=MatchLowercase}
  \defaultfontfeatures[\rmfamily]{Ligatures=TeX,Scale=1}
\fi
\usepackage{lmodern}
\ifPDFTeX\else
  % xetex/luatex font selection
\fi
% Use upquote if available, for straight quotes in verbatim environments
\IfFileExists{upquote.sty}{\usepackage{upquote}}{}
\IfFileExists{microtype.sty}{% use microtype if available
  \usepackage[]{microtype}
  \UseMicrotypeSet[protrusion]{basicmath} % disable protrusion for tt fonts
}{}
\makeatletter
\@ifundefined{KOMAClassName}{% if non-KOMA class
  \IfFileExists{parskip.sty}{%
    \usepackage{parskip}
  }{% else
    \setlength{\parindent}{0pt}
    \setlength{\parskip}{6pt plus 2pt minus 1pt}}
}{% if KOMA class
  \KOMAoptions{parskip=half}}
\makeatother
\setlength{\emergencystretch}{3em} % prevent overfull lines
\providecommand{\tightlist}{%
  \setlength{\itemsep}{0pt}\setlength{\parskip}{0pt}}
\usepackage{bookmark}
\IfFileExists{xurl.sty}{\usepackage{xurl}}{} % add URL line breaks if available
\urlstyle{same}
\hypersetup{
  hidelinks,
  pdfcreator={LaTeX via pandoc}}

\author{}
\date{}

\begin{document}

\section{2025-05-16 Introduction to GP}\label{introduction-to-gp}

\emph{Converted from PDF: 2025-05-16 Introduction to GP.pdf}

\begin{center}\rule{0.5\linewidth}{0.5pt}\end{center}

\subsection{Page 1}\label{page-1}

An Introduction to Gaussian Processes (GP) Le Yang, Joe Chen and Jeremy
Waston Department of ECE, University of Canterbury

\begin{center}\rule{0.5\linewidth}{0.5pt}\end{center}

\subsection{Page 2}\label{page-2}

Outline • Supervised learning and regression • Fundamentals of Gaussian
Processes (GP) • Mean function, covariance function and hyperparameters
in GP • GP for nonlinear regression • GP in spatiotemporal modeling •
References

\begin{center}\rule{0.5\linewidth}{0.5pt}\end{center}

\subsection{Page 3}\label{page-3}

Supervised Learning • Supervised learning • Establish (generally
non-linear) curve(s) to • Build decision boundaries (classification) •
Fit through labeled data (regression) Classification Problem Regression
Problem

\begin{center}\rule{0.5\linewidth}{0.5pt}\end{center}

\subsection{Page 4}\label{page-4}

Supervised Regression • Formulation of a regression problem • Given a
set of 𝑁 input-output pairs (labeled training data) \{𝒙𝑛, 𝑦𝑛\}𝑛=1 𝑁 •
𝒙𝑛: input • 𝑦𝑛: output • Find a nonlinear (multivariate) function 𝑦=
𝑓(𝒙) that can • Describe the training data well (no underfitting) •
Generalize to inputs other than 𝑥𝑛 well (no overfitting) •
Predictive/generative modeling

\begin{center}\rule{0.5\linewidth}{0.5pt}\end{center}

\subsection{Page 5}\label{page-5}

Fundamentals of Gaussian Processes (1) • Definition • A Gaussian
Processes (GP) is a collection of random variables with the joint
distribution of any subset being Gaussian • To specify a (multivariate)
Gaussian distribution, we only need to provide • Mean and covariance •
Similarly, to specify a GP, we need • Mean function • Covariance
function

\begin{center}\rule{0.5\linewidth}{0.5pt}\end{center}

\subsection{Page 6}\label{page-6}

Fundamentals of Gaussian Processes (2) • Impose GP prior on the function
𝑦= 𝑓(𝒙) to be estimated • Any subset of function values 𝑦𝑖= 𝑓(𝒙𝑖) is
Gaussian • 𝑝(𝑦1, 𝑦2, \ldots{} , 𝑦𝑀) is Gaussian • 𝑀= 1,2,3, \ldots{} is
an arbitrary positive integer • What are the mean and covariance for
𝑝(𝑦1, 𝑦2, \ldots{} , 𝑦𝑀)? • Note that the inputs for 𝑦𝑖 are different •
Cannot use fixed mean and covariance • Mean and covariance should depend
on the input 𝒙1, 𝒙2, \ldots{} , 𝒙𝑀

\begin{center}\rule{0.5\linewidth}{0.5pt}\end{center}

\subsection{Page 7}\label{page-7}

Fundamentals of Gaussian Processes (3) • Mean function 𝑚(𝒙) • Specify
the mean vector for 𝑝(𝑦1, 𝑦2, \ldots{} , 𝑦𝑀) • 𝝁𝑀×1 = 𝑚(𝒙1) ⋮ 𝑚(𝒙𝑀) •
Covariance function 𝑘𝒙, 𝒙′ • Specify the covariance matrix for 𝑝(𝑦1, 𝑦2,
\ldots{} , 𝑦𝑀) • 𝜮𝑀×𝑀= 𝑘(𝒙1, 𝒙1) ⋯ 𝑘(𝒙1, 𝒙𝑀) ⋮ ⋱ ⋮ 𝑘(𝒙𝑀, 𝒙1) ⋯ 𝑘(𝒙𝑀, 𝒙𝑀)

\begin{center}\rule{0.5\linewidth}{0.5pt}\end{center}

\subsection{Page 8}\label{page-8}

Mean Function and Covariance Function in GP (1) • Mean function 𝑚(𝒙)
captures the trend in the function 𝑓(𝒙) • 𝑚𝒙= 0 is commonly used and in
many cases, a good choice • 𝑚𝒙= σ𝑗=1 𝐽 𝑎𝑗𝜑𝑗(𝒙), 𝜑𝑗(𝒙) is a
linear/nonlinear function of 𝒙 (surrogate ☺) • Covariance function 𝑘𝒙,
𝒙′\\
• Specifies the variance of the function value when 𝒙= 𝒙′ • Specifies
the correlation between function values 𝑓(𝒙) and 𝑓(𝒙′) •
𝑓𝒙\textasciitilde 𝒢𝒫𝑚𝒙, 𝑘𝒙, 𝒙′ : 𝑓(𝒙) has a GP prior with mean function
𝑚𝒙 and covariance function 𝑘𝒙, 𝒙′

\begin{center}\rule{0.5\linewidth}{0.5pt}\end{center}

\subsection{Page 9}\label{page-9}

Mean Function and Covariance Function in GP (2) • Squared exponential
(SE): a commonly used covariance function • 𝑘𝒙, 𝒙′ = 𝜎𝑓 2exp −
\textbar\textbar 𝒙−𝒙′\textbar\textbar2 2𝑙2 • 𝑘𝒙, 𝒙= 𝜎𝑓 2 • 𝜎𝑓 2: a
priori variance of function values 𝑓(𝒙) • 𝑙 : characteristic length •
Smoothness of the function 𝑓(𝑥) to be modeled • What if the training
data has noise in 𝑦𝑛? • 𝑘𝒙, 𝒙′ = 𝜎𝑓 2 exp − 𝒙−𝒙′ 2 2𝑙2 + 𝜎𝑣2𝛿(𝒙−𝒙′)\\
• 𝜎𝑣2 : noise variance • 𝛿(∙) : Dirac delta function

\begin{center}\rule{0.5\linewidth}{0.5pt}\end{center}

\subsection{Page 10}\label{page-10}

Mean Function and Covariance Function in GP (3) • Sampling from
𝑧\textasciitilde 𝑁𝜇, Σ v.s. sampling from 𝑓𝑥\textasciitilde 𝒢𝒫𝑚𝑥, 𝑘𝑥, 𝑥′
𝜇= 0 0 , Σ = 1 0.9 0.9 1 𝑚𝑥= 0, 𝑘𝑥, 𝑥′ =exp −𝑥−𝑥′ 2 + 𝛿(𝑥−𝑥′)

\begin{center}\rule{0.5\linewidth}{0.5pt}\end{center}

\subsection{Page 11}\label{page-11}

Hyperparameters in GP (1) • Hyperparameters in GP • Unknown parameters
in the mean function 𝑚𝒙 and covariance function 𝑘𝒙, 𝒙′ • Example • 𝑚𝒙=
σ𝑗=1 𝐽 𝑎𝑗𝜑𝑗(𝒙) • 𝑘𝒙, 𝒙′ = 𝜎𝑓 2 exp − 𝒙−𝒙′ 2 2𝑙2 + 𝜎𝑣2𝛿(𝒙−𝒙′)\\
• 𝐽+ 3 hyperparameters: \{𝑎1, 𝑎2, \ldots{} , 𝑎𝐽, 𝑙, 𝜎𝑓 2, 𝜎𝑣2\}

\begin{center}\rule{0.5\linewidth}{0.5pt}\end{center}

\subsection{Page 12}\label{page-12}

Hyperparameters in GP (2) • Sampling from 𝑓𝑥\textasciitilde 𝒢𝒫𝑚𝑥, 𝑘𝑥, 𝑥′
• 𝑚𝑥= 0 • 𝑘𝑥, 𝑥′ = 𝜎𝑓 2 exp − 𝑥−𝑥′ 2 2𝑙2 + 𝜎𝑣2𝛿(𝑥−𝑥′) • 𝑙= 1 • 𝜎𝑣= 0 :
no measurement noise • (a) 𝜎𝑓= 1 • (b) 𝜎𝑓= 2 • (c) 𝜎𝑓= 3 • (d) 𝜎𝑓= 0.5 •
𝜎𝑓 specifies the `signal' energy

\begin{center}\rule{0.5\linewidth}{0.5pt}\end{center}

\subsection{Page 13}\label{page-13}

Hyperparameters in GP (3) • Sampling from 𝑓𝑥\textasciitilde 𝒢𝒫𝑚𝑥, 𝑘𝑥, 𝑥′
• 𝑚𝑥= 0 • 𝑘𝑥, 𝑥′ = 𝜎𝑓 2 exp − 𝑥−𝑥′ 2 2𝑙2 + 𝜎𝑣2𝛿(𝑥−𝑥′) • 𝜎𝑣= 0 : no
measurement noise • 𝜎𝑓= 1 • (a) 𝑙= 1 • (b) 𝑙= 3 • (c) 𝑙= 5 • (d) 𝑙= 0.5
• 𝑙 specifies the smoothness of the function

\begin{center}\rule{0.5\linewidth}{0.5pt}\end{center}

\subsection{Page 14}\label{page-14}

Hyperparameters in GP (4) • Sampling from 𝑓𝑥\textasciitilde 𝒢𝒫𝑚𝑥, 𝑘𝑥, 𝑥′
• 𝑚𝑥= 0 • 𝑘𝑥, 𝑥′ = 𝜎𝑓 2 exp − 𝑥−𝑥′ 2 2𝑙2 + 𝜎𝑣2𝛿(𝑥−𝑥′) • 𝑙= 1 • 𝜎𝑓=1 •
(a) 𝜎𝑣= 0 • (b) 𝜎𝑣= 0.01 • (c) 𝜎𝑣= 0.05 • (d) 𝜎𝑣= 0.1 • 𝜎𝑣 specifies the
SNR

\begin{center}\rule{0.5\linewidth}{0.5pt}\end{center}

\subsection{Page 15}\label{page-15}

Hyperparameters in GP (5) • Maximum likelihood (ML) estimation of GP
hyperparameters • 𝑓𝒙\textasciitilde 𝒢𝒫𝑚𝒙, 𝑘𝒙, 𝒙′ • The training data
satisfy 𝑝(𝑦1, 𝑦2, \ldots{} , 𝑦𝑁) = 𝑁 𝑚(𝒙1) ⋮ 𝑚(𝒙𝑁) , 𝑘(𝒙1, 𝒙1) ⋯ 𝑘(𝒙1,
𝒙𝑁) ⋮ ⋱ ⋮ 𝑘(𝒙𝑁, 𝒙1) ⋯ 𝑘(𝒙𝑁, 𝒙𝑁) • Find the hyperparameters through
solving via conjugate gradient 𝑚𝑎𝑥log(𝑝(𝑦1, 𝑦2, \ldots{} , 𝑦𝑁))

\begin{center}\rule{0.5\linewidth}{0.5pt}\end{center}

\subsection{Page 16}\label{page-16}

GP for Nonlinear Regression (1) • Problem setup • Given a set of 𝑁
training data points \{𝒙𝑛, 𝑦𝑛\}𝑛=1 𝑁 • The function to be modeled, 𝑓𝒙,
has a GP prior 𝑓𝒙\textasciitilde 𝒢𝒫𝑚𝒙, 𝑘𝒙, 𝒙′ • GP hyperparameters have
been found through e.g., ML estimation • Data-driven approach • We are
interested to find the posterior 𝑝(𝑦∗\textbar 𝒙∗, \{𝒙𝑛, 𝑦𝑛\}𝑛=1 𝑁 ) •
𝑦∗= 𝑓(𝒙∗): the function value at an input 𝒙∗ • `Predictive' probability
of 𝑦∗ • Modeling of mapping from 𝒙∗ to 𝑦∗

\begin{center}\rule{0.5\linewidth}{0.5pt}\end{center}

\subsection{Page 17}\label{page-17}

GP for Nonlinear Regression (2) • By Bayes theorem 𝑝(𝑦∗\textbar 𝒙∗,
\{𝒙𝑛, 𝑦𝑛\}𝑛=1 𝑁 ) = 𝑝( 𝒙∗, 𝑦∗, \{𝒙𝑛, 𝑦𝑛\}𝑛=1 𝑁 ) 𝑝(\{𝒙𝑛, 𝑦𝑛\}𝑛=1 𝑁 ) •
Note that because 𝑓𝒙\textasciitilde 𝒢𝒫𝑚𝒙, 𝑘𝒙, 𝒙′ , 𝑝 𝒙𝑛, 𝑦𝑛𝑛=1 𝑁 = 𝑝(𝑦1,
𝑦2, \ldots{} , 𝑦𝑁) = 𝑁 𝑚(𝒙1) ⋮ 𝑚(𝒙𝑁) , 𝑘(𝒙1, 𝒙1) ⋯ 𝑘(𝒙1, 𝒙𝑁) ⋮ ⋱ ⋮ 𝑘(𝒙𝑁,
𝒙1) ⋯ 𝑘(𝒙𝑁, 𝒙𝑁) 𝑝 𝒙∗, 𝑦∗, 𝒙𝑛, 𝑦𝑛𝑛=1 𝑁 = 𝑝(𝑦∗, 𝑦1, 𝑦2, \ldots{} , 𝑦𝑁) = 𝑁
𝑚(𝒙∗) 𝑚(𝒙1) ⋮ 𝑚(𝒙𝑁) , 𝑘(𝒙∗, 𝒙∗) 𝑘(𝒙∗, 𝒙1) ⋯ 𝑘(𝒙∗, 𝒙1) 𝑘(𝒙1, 𝒙∗) ⋮ 𝑘(𝒙𝑁,
𝒙∗) 𝑘(𝒙1, 𝒙1) ⋯ 𝑘(𝒙1, 𝒙𝑁) ⋮ ⋱ ⋮ 𝑘(𝒙𝑁, 𝒙1) ⋯ 𝑘(𝒙𝑁, 𝒙𝑁)

\begin{center}\rule{0.5\linewidth}{0.5pt}\end{center}

\subsection{Page 18}\label{page-18}

GP for Nonlinear Regression (3) • 𝑝(𝑦∗\textbar 𝒙∗, \{𝒙𝑛, 𝑦𝑛\}𝑛=1 𝑁 ) is
Gaussian • By Gaussian conditional, 𝑝(𝑦∗\textbar 𝒙∗, \{𝒙𝑛, 𝑦𝑛\}𝑛=1 𝑁 ) =
𝑁(𝜇∗, Σ∗) where 𝜇∗= 𝑚𝒙∗+ {[}𝑘(𝒙∗, 𝒙1) ⋯ 𝑘(𝒙∗, 𝒙𝑁){]} ∙ 𝑘(𝒙1, 𝒙1) ⋯ 𝑘(𝒙1,
𝒙𝑁) ⋮ ⋱ ⋮ 𝑘(𝒙𝑁, 𝒙1) ⋯ 𝑘(𝒙𝑁, 𝒙𝑀) −1 𝑦1 ⋮ 𝑦𝑁 − 𝑚(𝒙1) ⋮ 𝑚(𝒙𝑁) Σ∗= 𝑘𝒙∗,
𝒙∗−{[}𝑘(𝒙∗, 𝒙1) ⋯ 𝑘(𝒙∗, 𝒙𝑁){]} ∙ 𝑘(𝒙1, 𝒙1) ⋯ 𝑘(𝒙1, 𝒙𝑁) ⋮ ⋱ ⋮ 𝑘(𝒙𝑁, 𝒙1) ⋯
𝑘(𝒙𝑁, 𝒙𝑀) −1 𝑘(𝒙1, 𝒙∗) ⋮ 𝑘(𝒙𝑁, 𝒙∗) • 𝜇∗: predictive function value at 𝒙∗
• Σ∗: prediction uncertainty

\begin{center}\rule{0.5\linewidth}{0.5pt}\end{center}

\subsection{Page 19}\label{page-19}

GP for Nonlinear Regression (3) • Example • 𝑚𝑥= 0 • 𝑘𝑥, 𝑥′ = 𝜎𝑓 2 exp −
𝑥−𝑥′ 2 2𝑙2 + 𝜎𝑣2𝛿(𝑥−𝑥′) • 𝑙= 1 • 𝜎𝑣= 0 : no measurement noise • 𝜎𝑓= 1 •
`Accurate' prediction around training data points • `Large' uncertainty
when away from training data

\begin{center}\rule{0.5\linewidth}{0.5pt}\end{center}

\subsection{Page 20}\label{page-20}

GP for Nonlinear Regression (4) • GP-based nonlinear regression is
non-parametric • Complexity proportional to 𝑂(𝑁3), scalable methods
needed for large-scale dataset • GP-based nonlinear regression
explicitly quantifies prediction uncertainty • GP-based nonlinear
regression is equivalent to a neural network with one hidden layer and
infinitely large number of neurons • Powerful function approximator • A
few mean functions and covariance functions are available for modeling •
Combinations of different mean functions and covariance functions
possible

\begin{center}\rule{0.5\linewidth}{0.5pt}\end{center}

\subsection{Page 21}\label{page-21}

GP for Spatiotemporal Modeling (1) • Spatiotemporal phenomenon is an
event relating to both time and space • El Nino dataset: 178,080
meteorological measurements across the Pacific • Air temperature •
Relative humidity • Surface winds • Sea surface temperature • 25 sensors
and span of 5 years (daily records) • Air temperature at one particular
location →

\begin{center}\rule{0.5\linewidth}{0.5pt}\end{center}

\subsection{Page 22}\label{page-22}

GP for Spatiotemporal Modeling (2) • Prediction at 40 uniformly spaced
time points using 30 measurements • Mean function 𝑚(𝑥) = 27 • SE
covariance function with 𝑙= 100 days, 𝜎𝑓= 1, 𝜎𝑣= 1 • Prediction at 200
uniformly spaced time points using 100 measurements • Mean function 𝑚(𝑥)
= 27 • SE covariance function with 𝑙= 100 days, 𝜎𝑓= 1, 𝜎𝑣= 1

\begin{center}\rule{0.5\linewidth}{0.5pt}\end{center}

\subsection{Page 23}\label{page-23}

References Book Chapters {[}1{]}. C. E. Rasmussen and C. K. I. Williams,
Gaussian Processes for Machine Learning, The MIT Press, 2006. {[}2{]}.
K. Murphy, Probabilistic Machine Learning: An introduction, Chapter 17,
Book Draft, 2021. Scalable GP: {[}3{]}. Liu et. al., When Gaussian
Processes meets big data: A review with scalable GPs, IEEE Trans. Neural
Networks and Learning Systems, 2020 Spatiotemporal Modeling: {[}4{]}. C.
Jidling, Tailoring Gaussian Processes for Tomographic Reconstruction,
PhD Dissertation, Uppsala University, 2019. {[}5{]}. Singh et. al.,
Efficient informative sensing using multiple robots, Journal of
Artificial Intelligence Research, 2009. {[}6{]}. Kopp et. al.,
Temperature-driven global sea-level variability in the common era, PNAS,
2016. {[}7{]}. Guestrin et. al., Near-optimal sensor placements in
Gaussian Processes. ICML, 2005.

\begin{center}\rule{0.5\linewidth}{0.5pt}\end{center}

\end{document}
